\section{Empirical Strategy}

\subsection{Experimental Analysis: Identification and Tests of H1–H3}

The first part of the empirical analysis examines how the experimentally varied characteristics of the migrant profiles 
affect respondents' beliefs about migrants' labor-market and social integration outcomes. The unit of observation is a 
vignette evaluation, where respondent \(i\) evaluates migrant profile \(j\). Since each respondent evaluates several 
migrant profiles, the analysis contains multiple observations for each respondent.

The factorial vignette design independently randomizes five characteristics across profiles: country of origin, reason 
for migration, gender, human capital, and German-language proficiency (see Chapter 3). Random assignment makes these 
characteristics independent of respondents’ pre-existing characteristics and, in expectation, of the other randomized 
profile attributes. The experimental variation can therefore be used to identify the causal effect of each presented 
migrant characteristic on respondents’ stated beliefs. In other words, the experiment identifies how respondents’ 
beliefs change when one characteristic of the migrant profile changes. For example, it identifies how respondents’ 
beliefs differ when an otherwise identical migrant is presented as Ukrainian rather than Syrian.

\subsubsection{Variables and Baseline Specification}

\noindent\textbf{Baseline specification.} The five outcome variables introduced in Section 3.3 are analyzed separately. 
Let \(Y_{ij}^{k}\) denote respondent \(i\)'s belief about outcome \(k\) for vignette profile \(j\), where
$$ k\in\{\mathrm{Emp},\mathrm{Inc},\mathrm{Hire},\mathrm{Work},\mathrm{Soc}\}. $$
\(Y_{ij}^{\mathrm{Emp}}\) denotes the expected probability of obtaining regular paid employment within two years, 
\(Y_{ij}^{\mathrm{Inc}}\) denotes expected monthly net income, \(Y_{ij}^{\mathrm{Hire}}\) and \(Y_{ij}^{\mathrm{Work}}\) 
denote expected discrimination in hiring and in the workplace, respectively, and \(Y_{ij}^{\mathrm{Soc}}\) denotes 
expected social integration.

The randomized vignette characteristics are represented by five binary variables. \(Ukr_j\) equals one if profile \(j\) 
describes a Ukrainian migrant and zero if the migrant is Syrian. \(Ref_j\) equals one if the migrant is presented as a 
war refugee and zero if the migrant is presented as an economic migrant. \(Fem_j\) equals one for a female migrant and 
zero for a male migrant. \(HighSkill_j\) equals one if the migrant has completed a university degree and zero if the 
migrant has no formal vocational qualification. Finally, \(German_j\) equals one if the migrant speaks German well and 
zero if the migrant speaks German only to a limited extent. These definitions imply that the reference profile 
is a male Syrian economic migrant without a formal vocational qualification who speaks German only to a limited extent.

For each outcome \(k\), the baseline specification is
\begin{equation}
\begin{aligned}
Y_{ij}^{k}
&= \alpha^{k}
+ \beta_{1}^{k} Ukr_{j}
+ \beta_{2}^{k} Ref_{j}
+ \beta_{3}^{k} Fem_{j}
+ \beta_{4}^{k} HighSkill_{j}
+ \beta_{5}^{k} German_{j}
+ \varepsilon_{ij}^{k}.
\end{aligned}
\label{eq:baseline}
\end{equation}
The coefficients \(\beta_1^k\) to \(\beta_5^k\) capture the average causal effects of the respective randomized vignette 
characteristics on outcome \(k\), relative to their reference categories and averaging over the randomized values of 
the remaining profile characteristics. The baseline specifications are estimated using linear regression, and the 
interpretation of the coefficients depends on the respective outcome measure. For the employment outcome, which is 
measured on a 0–100 probability scale, coefficients are expressed in percentage points. For the income outcome, the 
seven ordered categories are coded from 1 to 7, so coefficients are expressed in income-category points rather than as 
specific differences in euros. For the remaining outcomes, which are measured on seven-point response scales, 
coefficients are expressed in scale points. Since the ordinal response categories cannot necessarily be assumed to 
be equally spaced, ordered logit models are additionally estimated as robustness checks.
The term \(\varepsilon_{ij}^{k}\) denotes the error term and captures variation in respondent \(i\)'s evaluation of 
profile \(j\) for outcome \(k\) that is not explained by the variables included in the regression.\bigskip


\noindent\textbf{Adjusted specification.} Because the vignette attributes are independently randomized, respondent-level controls are not required for 
identification of the coefficients in equation (1). The unadjusted specification therefore constitutes the main 
specification for estimating the experimental effects. As an adjusted specification, respondent background 
characteristics can additionally be included:
\begin{equation}
\begin{aligned}
Y_{ij}^{k}
&= \alpha^{k}
+ \beta_{1}^{k}Ukr_j
+ \beta_{2}^{k}Ref_j
+ \beta_{3}^{k}Fem_j
+ \beta_{4}^{k}HighSkill_j \\
&\quad
+ \beta_{5}^{k}German_j
+ \boldsymbol{\delta}^{k\prime}X_i
+ \varepsilon_{ij}^{k}.
\end{aligned}
\end{equation}
where \(X_i\) is a vector of respondent-level controls including age, gender, education, employment status, household 
net income, migration background, region of residence in Germany, religion, and political self-placement, and is denoted by
\[
X_i =
\left(
Age_i,\,
Gender_i,\,
Edu_i,\,
Emp_i,\,
Inc_i,\,
Mig_i,\,
East_i,\,
Religion_i,\,
Right_i
\right)'.
\]
The inclusion of \(X_i\) does not provide the source of identification for the randomized vignette effects, but this
adjusted specification can be used to assess the robustness and potentially improve the precision of the experimental 
estimates. The construction and use of the respondent-level controls are discussed further in Section 4.4.\bigskip

\noindent\textbf{Standard errors.} Since each respondent evaluates multiple vignette profiles, observations belonging to the same respondent cannot be 
treated as statistically independent. Standard errors are therefore clustered at the respondent level in all 
vignette-level specifications, allowing the regression residuals to be correlated across the profiles evaluated by the 
same respondent. This accounts for the repeated-evaluation structure of the proposed experiment.

\subsubsection{Main Effects: H1a, H1b, H2, H3a, and H3b}

Since several hypotheses include more than one outcome variable, equation (1) is estimated separately for each relevant 
outcome. For each hypothesis, the coefficient of interest and its predicted sign are specified below. The estimated 
coefficient shows the direction and size of the effect, while 95\% confidence intervals are reported to indicate the 
uncertainty around the estimate. The estimated effect is reported and interpreted separately for each outcome. For 
example, if the results differ between employment and income, this difference is reported explicitly rather than 
combining the outcomes into a single measure of support for the hypothesis. Whether the results provide support for 
the respective hypothesis is determined from the pattern of the estimates, considering their predicted 
direction, magnitude, and statistical uncertainty. This avoids imposing a strict rule for hypothesis support before the empirical results are observed.
Since testing multiple outcomes also raises a multiple-testing concern, this is accounted for in the statistical 
inference, which will further be discussed in Section 4.4. 
Equation (2), which additionally includes respondent-level controls, is estimated as a robustness check.\bigskip

\noindent\textbf{Testing H1a.} H1a predicts that migrants with higher human capital are expected to achieve more 
favorable labor-market integration outcomes. Equation (1) is therefore estimated for the employment and income 
outcomes, \(Y_{ij}^{\mathrm{Emp}}\) and \(Y_{ij}^{\mathrm{Inc}}\). The coefficient of interest is \(\beta_{4}^{k}\) on \(HighSkill_j\), 
which captures the causal effect of presenting a migrant with a university degree rather than with no formal vocational 
qualification on respondents’ beliefs about the respective outcome. H1a predicts
$$ \beta_{4}^{\mathrm{Emp}}>0 \qquad \text{and} \qquad \beta_{4}^{\mathrm{Inc}}>0. $$
Thus, positive coefficients indicate more favorable beliefs about the respective labor-market outcomes for profiles 
with higher human capital.\bigskip

\noindent\textbf{Testing H1b.} H1b predicts that migrants with higher German-language proficiency are expected to 
achieve more favorable labor-market and social-integration outcomes. Equation (1) is therefore estimated
\(Y_{ij}^{\mathrm{Emp}}\), \(Y_{ij}^{\mathrm{Inc}}\), and expected social integration \(Y_{ij}^{\mathrm{Soc}}\). 
The coefficient of interest is \(\beta_{5}^{k}\) on \(German_j\), which captures the causal effect of presenting a 
migrant as speaking German well rather than only to a limited extent on respondents’ beliefs. H1b predicts
$$ \beta_{5}^{\mathrm{Emp}}>0, \qquad \beta_{5}^{\mathrm{Inc}}>0, \qquad \text{and} \qquad \beta_{5}^{\mathrm{Soc}}>0. $$
Thus, positive coefficients indicate more favorable beliefs about the respective outcomes for profiles with higher 
German-language proficiency.\bigskip

\noindent\textbf{Testing H2.} H2 predicts that refugees are expected to achieve less favorable labor-market integration 
outcomes than otherwise comparable economic migrants. Equation (1) is estimated for \(Y_{ij}^{\mathrm{Emp}}\) and 
\(Y_{ij}^{\mathrm{Inc}}\). The coefficient of interest is \(\beta_{2}^{k}\) on \(Ref_j\), which captures the causal effect of 
presenting a migrant as a war refugee rather than as an economic migrant on respondents’ beliefs. H2 predicts
$$ \beta_{2}^{\mathrm{Emp}}<0 \qquad \text{and} \qquad \beta_{2}^{\mathrm{Inc}}<0. $$
Thus, negative coefficients indicate less favorable beliefs about the respective outcomes for profiles 
presented as refugees rather than economic migrants.\bigskip

\noindent\textbf{Testing H3a.} H3a predicts that Ukrainian migrants are expected to achieve more favorable labor-market 
and social-integration outcomes than otherwise comparable Syrian migrants. Equation (1) is estimated for 
\(Y_{ij}^{\mathrm{Emp}}\), \(Y_{ij}^{\mathrm{Inc}}\), and \(Y_{ij}^{\mathrm{Soc}}\). The coefficient of interest is 
\(\beta_{1}^{k}\) on \(Ukr_j\), which captures the causal effect of presenting a migrant as Ukrainian rather than 
Syrian on respondents’ beliefs. H3a predicts
$$ \beta_{1}^{\mathrm{Emp}}>0, \qquad \beta_{1}^{\mathrm{Inc}}>0, \qquad \text{and} \qquad \beta_{1}^{\mathrm{Soc}}>0. $$
Thus, positive coefficients indicate more favorable beliefs about the respective outcomes for profiles presented as 
Ukrainian rather than Syrian.\bigskip

\noindent\textbf{Testing H3b.} H3b examines country-of-origin differences in expected labor-market discrimination and 
predicts that Syrian migrants are expected to face greater discrimination than otherwise comparable Ukrainian migrants. 
The hypothesis is tested using the country-of-origin indicator, \(Ukr_j\), with \(Y_{ij}^{\mathrm{Hire}}\) and 
\(Y_{ij}^{\mathrm{Work}}\) as the dependent variables. An important difference from the preceding outcomes is that 
higher values of both discrimination measures indicate a less favorable outcome, namely greater expected discrimination. 
Since Syria is the reference category and \(Ukr_j=1\) represents Ukraine, H3b therefore predicts
$$ \beta_{1}^{\mathrm{Hire}}<0 \qquad \text{and} \qquad \beta_{1}^{\mathrm{Work}}<0. $$
Negative coefficients thus indicate that profiles presented as Ukrainian are expected to face less hiring and workplace 
discrimination than otherwise comparable Syrian profiles.

\subsubsection{Interaction Effects: H3c and Additional Analysis}

\noindent\textbf{Testing H3c.} While H3a and H3b examine average country-of-origin differences in expected integration 
and discrimination outcomes, H3c examines whether the size of these differences varies systematically between the two reasons for migration. 
Specifically, H3c predicts that differences between Ukrainian and Syrian migrants are greater when migrants are 
presented as refugees than when they are presented as economic migrants. 

To test this hypothesis, an interaction term between country of origin and reason for migration is added to the baseline specification:
\begin{equation}
\begin{aligned}
Y_{ij}^{k}
&= \alpha^{k}
+ \beta_{1}^{k}Ukr_j
+ \beta_{2}^{k}Ref_j
+ \beta_{3}^{k}Fem_j
+ \beta_{4}^{k}HighSkill_j
+ \beta_{5}^{k}German_j \\
&\quad
+ \beta_{6}^{k}(Ukr_j \times Ref_j)
+ \varepsilon_{ij}^{k}.
\end{aligned}
\end{equation}
Since \(Ref_j=0\) represents economic migrants, \(\beta_{1}^{k}\) captures the Ukrainian–Syrian difference among 
economic migrants. When \(Ref_j=1\), the corresponding Ukrainian–Syrian difference among refugees is given by 
\(\beta_{1}^{k}+\beta_{6}^{k}\). Thus, \(\beta_{6}^{k}\) captures how the country-of-origin difference changes 
when migrants are presented as refugees rather than economic migrants. For employment, income, and social integration, 
higher values represent more favorable outcomes and H3c therefore predicts
$$ \beta_{6}^{\mathrm{Emp}}>0,\qquad \beta_{6}^{\mathrm{Inc}}>0,\qquad \beta_{6}^{\mathrm{Soc}}>0. $$
For hiring and workplace discrimination, higher values instead represent greater expected discrimination. 
Since Ukrainian migrants are expected to face less discrimination, a larger country-of-origin difference among refugees 
corresponds to a more negative Ukrainian–Syrian difference. H3c therefore predicts
$$ \beta_{6}^{\mathrm{Hire}}<0 \qquad\text{and}\qquad \beta_{6}^{\mathrm{Work}}<0. $$

In addition to the interaction coefficient, the conditional country-of-origin difference among refugees, 
\(\beta_{1}^{k}+\beta_{6}^{k}\), is reported together with its 95\% confidence interval. This shows whether the 
Ukrainian–Syrian difference among refugees itself is in the direction predicted by H3a and H3b, while 
\(\beta_{6}^{k}\) tests whether this difference is greater among refugees than among economic migrants.\bigskip

\noindent\textbf{Additional human-capital analysis.} In addition to H3c, the analysis similarly explores whether the 
Ukrainian–Syrian difference in expected integration outcomes varies with migrants’ human capital. For this purpose, 
the equation (1) is extended separately by the interaction \(Ukr_j\times HighSkill_j\). If higher human capital reduces 
the Ukrainian–Syrian gap in expected integration outcomes, the coefficient on this interaction is expected to be 
negative. This analysis is treated as supporting rather than as a separate hypothesis.

\subsection{Associational Analysis: Contact, Experience, Media, and Memory}

Section 4.2 examines H4–H6, which concern the relationship between respondents’ personal contact, 
media exposure, most memorable experiences and information and their beliefs about migrant integration. 
Unlike the vignette characteristics analyzed in Section 4.1, these respondent-level measures are not randomly 
assigned. The corresponding coefficients are therefore interpreted as conditional associations rather than 
causal effects. The analysis controls for relevant observed characteristics and, where appropriate, exploits 
within-respondent differences between experiences with Ukrainian and Syrian migrants. These specifications can 
reduce some sources of confounding but cannot rule out selection, reverse causality, or unobserved factors that 
may jointly influence respondents’ experiences, information exposure, and beliefs. These limitations are considered 
further in Section 4.4. 

\subsubsection{Contact and Experience: H4a–H4b}

\noindent\textbf{Testing H4a.} H4a predicts that, conditional on contact frequency, more positive overall personal 
experiences with a migrant group are associated with more favorable beliefs about its integration outcomes.

Because contact and experience are measured separately for Ukrainian and Syrian migrants, these measures are matched 
to the country of origin presented in each vignette. Let \(g(j)\in\{\mathrm{Ukrainian},\mathrm{Syrian}\}\) denote the 
country-of-origin group presented in vignette \(j\). Thus, \(C_{i,g(j)}\) and \(Q_{i,g(j)}\) refer to respondent \(i\)'s 
contact frequency and overall contact quality, respectively, with the group presented in vignette \(j\). 
For example, a Ukrainian vignette is matched to the respondent's reported contact and experience with Ukrainian migrants.

To test H4a, the following specification is estimated:
\begin{equation}
Y_{ij}^{k}
=
\alpha^{k}
+ \boldsymbol{\beta}^{k\prime} Z_j
+ \gamma_{C}^{k} C_{i,g(j)}
+ \gamma_{Q}^{k} Q_{i,g(j)}
+ \boldsymbol{\delta}^{k\prime} X_i
+ \varepsilon_{ij}^{k},
\end{equation}
where \(Z_j\) is the vector of randomized vignette characteristics,
\[
Z_j =
\left(
Ukr_j,\,
Ref_j,\,
Fem_j,\,
HighSkill_j,\,
German_j
\right)'.
\]
\(X_i\) contains the respondent-level controls defined above. The coefficient of interest for H4a is \(\gamma_Q^k\), 
which captures the association between the overall quality of contact with the relevant migrant group and respondent 
beliefs. \(C_{i,g(j)}\), is included so that the association between overall contact quality and beliefs is estimated 
conditional on how frequently respondents have contact with the respective migrant group. Its coefficient, 
\(\gamma_C^k\), captures the association between contact frequency and beliefs, holding overall contact quality
and the other included variables constant.

H4a predicts that more positive contact is associated with more favorable beliefs about employment, income, and 
social-integration outcomes:
\[
\gamma_Q^{\mathrm{Emp}}>0,\qquad
\gamma_Q^{\mathrm{Inc}}>0,\qquad
\gamma_Q^{\mathrm{Soc}}>0.
\]

\noindent\textbf{Testing H4b.} H4b examines whether negative personal experiences are more strongly associated with 
respondents' beliefs than positive personal experiences. Using the positive and negative experience measures 
\(Positive_{ig}\) and \(Negative_{ig}\) constructed from \(Q_{ig}\) as described in Section 3.4.1, the two measures are 
included simultaneously in the following specification:
\begin{equation}
Y_{ij}^{k}
=
\alpha^{k}
+ \boldsymbol{\beta}^{k\prime}Z_j
+ \gamma_C^{k}C_{i,g(j)}
+ \gamma_P^{k}Positive_{i,g(j)}
+ \gamma_N^{k}Negative_{i,g(j)}
+ \boldsymbol{\delta}^{k\prime}X_i
+ \varepsilon_{ij}^{k}.
\end{equation}
For the employment, income, and social-integration outcomes, more positive experiences are expected to be associated 
with more favorable beliefs, whereas more negative experiences are expected to be associated with less favorable beliefs. 
Thus,
\[
\gamma_P^{\mathrm{Emp}},\,
\gamma_P^{\mathrm{Inc}},\,
\gamma_P^{\mathrm{Soc}}>0,
\qquad
\gamma_N^{\mathrm{Emp}},\,
\gamma_N^{\mathrm{Inc}},\,
\gamma_N^{\mathrm{Soc}}<0.
\]
H4b, however, predicts not only opposite directions but also a stronger association for negative than for equally 
strong positive experiences. This prediction can be tested by comparing the absolute magnitudes of their coefficients. 
For each relevant outcome \(k\), H4b predicts
\[
-\gamma_N^{k}>\gamma_P^{k},
\]
or equivalently,
\[
\gamma_P^{k}+\gamma_N^{k}<0.
\]
The comparison therefore tests whether moving one unit toward a more negative experience is associated with a larger 
change in beliefs than an equivalent one-unit movement toward a more positive experience.

\subsubsection{Media Exposure and Tone: H5}

\noindent\textbf{Testing H5.} H5 examines whether the overall valence of media information concerning a migrant group 
is associated with respondents' beliefs about that group's integration outcomes, conditional on the frequency of media 
exposure. As in the contact analysis, the group-specific media measures are matched to the country of origin presented 
in each vignette. The following specification is estimated:
\begin{equation}
Y_{ij}^{k}
=
\alpha^{k}
+ \boldsymbol{\beta}^{k\prime}Z_j
+ \gamma_M^{k}M_{i,g(j)}
+ \gamma_T^{k}T_{i,g(j)}
+ \boldsymbol{\delta}^{k\prime}X_i
+ \varepsilon_{ij}^{k}.
\end{equation}
The coefficient of interest is \(\gamma_T^{k}\) on \(T_{i,g(j)}\), which captures the association between the overall 
tone of media information and respondent beliefs. \(M_{i,g(j)}\) is included so that the association between media tone 
and beliefs is estimated conditional on how frequently respondents encounter media information about the respective 
migrant group. Its coefficient, \(\gamma_M^k\), captures the association between media-exposure frequency and beliefs, 
holding media tone and the other included variables constant.

For the employment, income, and social-integration outcomes, H5 predicts that more positive media information is 
associated with more favorable beliefs:
\[
\gamma_T^{\mathrm{Emp}}>0,\qquad
\gamma_T^{\mathrm{Inc}}>0,\qquad
\gamma_T^{\mathrm{Soc}}>0.
\]

\subsubsection{Memorable Experiences and Information: H6a–H6b}

\noindent\textbf{Testing H6a.} H6a examines whether the valence of a respondent's most memorable personal experience 
with a migrant group is associated with beliefs, conditional on contact frequency and overall contact quality. 
The valence of the most memorable experience, \(V_{ig}^{\mathrm{E}}\), is observed only when the respondent reports 
that such an experience comes to mind (\(E_{ig}=1\)). The following specification is therefore estimated only for the 
observations satisfying \(E_{i,g(j)}=1\):
\begin{equation}
Y_{ij}^{k}
=
\alpha^{k}
+ \boldsymbol{\beta}^{k\prime}Z_j
+ \gamma_C^{k}C_{i,g(j)}
+ \gamma_Q^{k}Q_{i,g(j)}
+ \gamma_E^{k}V_{i,g(j)}^{\mathrm{E}}
+ \boldsymbol{\delta}^{k\prime}X_i
+ \varepsilon_{ij}^{k}.
\end{equation}
The coefficient of interest is \(\gamma_E^{k}\). Thus, H6a examines whether memorable-experience valence is associated 
with beliefs even after respondents' overall experiences with the respective migrant group are taken into account. 
The direction, magnitude, and statistical uncertainty of \(\gamma_E^{k}\) are reported for each relevant outcome.\bigskip

\noindent\textbf{Testing H6b.} H6b examines the corresponding relationship for respondents' most memorable media 
information. Similarly, \(V_{ig}^{\mathrm{I}}\) is observed only when the respondent reports that a particularly 
memorable piece of media information exists, \(I_{ig}=1\). Then, the following specification is estimated:
\begin{equation}
Y_{ij}^{k}
=
\alpha^{k}
+ \boldsymbol{\beta}^{k\prime}Z_j
+ \gamma_M^{k}M_{i,g(j)}
+ \gamma_T^{k}T_{i,g(j)}
+ \gamma_I^{k}V_{i,g(j)}^{\mathrm{I}}
+ \boldsymbol{\delta}^{k\prime}X_i
+ \varepsilon_{ij}^{k}.
\end{equation}
The coefficient of interest is \(\gamma_I^{k}\). H6b therefore examines whether the valence of memorable media 
information is associated with beliefs conditional on respondents’ overall assessment of the media information they 
encounter. The direction, magnitude, and statistical uncertainty of \(\gamma_I^{k}\) are reported for each relevant 
outcome.\bigskip

\noindent\textbf{Supporting analyses.} In addition to H6a and H6b, I propose comparing the coefficients of the overall 
and memorable measures to examine whether their associations with respondents’ beliefs differ. This suggests comparing 
\(\gamma_Q^{k}\) and \(\gamma_E^{k}\) for contact experiences, and \(\gamma_T^{k}\) and \(\gamma_I^{k}\) for media 
information. Since the measures are recorded on the same seven-point scales, this makes it possible to test whether 
the corresponding coefficients differ from one another. The comparisons are interpreted cautiously because respondents’ 
overall assessments may themselves partly reflect their most memorable experiences or information. Conversely, the 
overall and memorable measures may also differ substantially, for example when a respondent reports generally positive 
experiences but recalls a particularly negative experience as most memorable. In such cases, the analysis can show 
whether beliefs are more strongly associated with the overall or the memorable measure.

As an additional exploratory analysis, the specifications used to examine H4--H6 are also estimated using the two 
discrimination outcomes, \(Y_{ij}^{\mathrm{Hire}}\) and \(Y_{ij}^{\mathrm{Work}}\). While H4--H6 focus on beliefs 
about integration outcomes, this extension examines whether respondents' contact experiences and media 
information are also associated with beliefs about the discrimination faced by the respective migrant group. 
The results are interpreted as exploratory associations.

\subsection{General Beliefs and Exploratory Analyses}

\subsection{Robustness, Data Treatment, and Interpretation}

explain the robustness of 1-7 and the hypothesis outcomes/decisions 
multiple testing problem

Citing measurement sources for the respondent questionnaire:
$C_{ig}^{Work}$ workplace contact frequency, 
$C_{ig}^{Social}$ social contact frequency,
$Q_{ig}$ contact quality.

Beliefs about large communities and network effects:

Large community
$\rightarrow$
More information/referrals
$\rightarrow$
Better employment

versus 

Large community
$\rightarrow$
Less native interaction
$\rightarrow$
Slower integration.

Construct:

$\Delta B_i^{\mathrm{Network}} = B_{i,\mathrm{S}}^{\mathrm{Network}} - B_{i,\mathrm{U}}^{\mathrm{Network}}$